% =================
% CADERNO DE PROVAS ICONOGRÁFICAS — design R1076
% =================
\chapter*{Caderno de Provas Iconográficas}
\addcontentsline{toc}{chapter}{Caderno de Provas Iconográficas}
\markboth{Caderno de Provas Iconográficas}{}

\noindent\textit{As imagens deste caderno pertencem à ficção documental do romance.
Não são reproduções de documentos históricos, nem provas de arquivo. Foram
compostas para este volume; o horror real de Barbacena e da seca não é
reconstituído aqui.}

\vspace{0.6em}
\noindent{\color{epi}\rule{0.22\textwidth}{1pt}}

\vspace{1em}
\begin{figure}[H]
  \centering
  \begin{minipage}[t]{0.46\textwidth}
    \centering
    \includegraphics[width=\linewidth,height=2.35in,keepaspectratio]{figures/immersao/placa_00.png}
    \caption*{\footnotesize\itshape Paciente 1.260. Retrato de arquivo.\\Reconstituição ficcional.}
  \end{minipage}\hfill
  \begin{minipage}[t]{0.46\textwidth}
    \centering
    \includegraphics[width=\linewidth,height=2.35in,keepaspectratio]{figures/immersao/placa_01.png}
    \caption*{\footnotesize\itshape Registro de admissão. Colônia de Barbacena, 1917.\\Reconstituição ficcional.}
  \end{minipage}
\end{figure}

\vspace{0.8em}
\begin{figure}[H]
  \centering
  \begin{minipage}[t]{0.46\textwidth}
    \centering
    \includegraphics[width=\linewidth,height=2.35in,keepaspectratio]{figures/immersao/placa_02.png}
    \caption*{\footnotesize\itshape Dr. Heitor Oliveira. Arquivo pessoal, 1979.\\Reconstituição ficcional.}
  \end{minipage}\hfill
  \begin{minipage}[t]{0.46\textwidth}
    \centering
    \includegraphics[width=\linewidth,height=2.35in,keepaspectratio]{figures/immersao/placa_03.png}
    \caption*{\footnotesize\itshape O homem do pano preto. Imagem não catalogada.\\Reconstituição ficcional.}
  \end{minipage}
\end{figure}

\vspace{0.8em}
\begin{figure}[H]
  \centering
  \begin{minipage}[t]{0.46\textwidth}
    \centering
    \includegraphics[width=\linewidth,height=2.35in,keepaspectratio]{figures/immersao/placa_04.png}
    \caption*{\footnotesize\itshape Dra. Lúcia Mendes. Rodoviária de Barbacena, 2026.\\Reconstituição ficcional.}
  \end{minipage}\hfill
  \begin{minipage}[t]{0.46\textwidth}
    \centering
    \includegraphics[width=\linewidth,height=2.35in,keepaspectratio]{figures/immersao/placa_05.png}
    \caption*{\footnotesize\itshape Fotografia 5. Arquivo pessoal de Lúcia Mendes, 2026.\\Reconstituição ficcional.}
  \end{minipage}
\end{figure}

\vspace{1em}
\noindent\textit{O Caderno não acrescenta provas: ele mostra o que a leitura já
viu. As imagens são parte do pacto ficcional, não do processo histórico.}

\clearpage
